\begin{center}\emph{Formen mit negativer Determinante.}\end{center}

\Article{171}

\textsc{Aufgabe.}\quad \emph{Zu einer beliebigen Form $(a,b,a')$ mit negativer Determinante $=-D$, wobei $D$ eine positive Zahl bezeichnet, eine eigentlich äquivalente Form $(A,B,C)$ finden, in der $A$ weder größer als $\surd\tfrac43D,$ $C$ noch kleiner als $2B$ ist.}

\emph{Lösung.}\quad Wir nehmen an, dass die gegebene Form nicht alle drei Bedingungen zugleich erfüllt; sonst müsste man keine andere Form suchen.\quad Sei $b'$ der absolut kleinste Rest der Zahl $-b$ modulo $a'$\HistoricalNote{*)}{Man beachte: Ist bei einer Form $(a,b,a')$ der erste oder letzte Term $a$ oder $a'$ $=0$, so ist ihre Determinante ein positives Quadrat. Das kann also im vorliegenden Fall nicht eintreten. --- Aus einem ähnlichen Grund können die äußeren Terme $a,a'$ einer Form mit negativer Determinante keine entgegengesetzten Vorzeichen haben.}, und sei $a''=\frac{b'b'+D}{a'}$. Dies ist eine ganze Zahl, weil $b'b'\equiv bb$, $b'b'+D\equiv bb+D\equiv aa'\equiv$ $0(\operatorname{mod.}a').$\quad Ist nun $a''<a'$, so sei wieder $b''$ der absolut kleinste Rest von $-b'$ modulo $a''$ und $a'''=\frac{b''b''+D}{a''}$. Ist auch hier $a'''<a''$, so sei erneut $b'''$ der absolut kleinste Rest von $-b''$ modulo $a'''$ und $a''''=\frac{b'''b'''+D}{a'''}$.\quad Man setzt dies fort, bis man in der Folge $a',a''$, $a''',a''''$ usw. einen Term $a^{m+1}$ erreicht, der nicht kleiner als sein Vorgänger $a^m$ ist. Das muss schließlich eintreten, da man sonst eine unendliche Folge ganzer Zahlen hätte, die immer weiter abnehmen.\quad Dann erfüllt die Form $(a^m,b^m,a^{m+1})$ alle Bedingungen.

\emph{Beweis.}\quad I.\quad In der Folge der Formen $(a,b,a'),\ (a',b',a''),\ (a'',b'',a''')$ usw. ist jede ihrer Vorgängerin benachbart; die letzte ist also der ersten eigentlich äquivalent (Artikel $159$, $160).$
\SourceLeaf{152}
II. Da $b^m$ der absolut kleinste Rest von $-b^{m-1}$ modulo $a^m$ ist, ist er nicht größer als $\frac12 a^m$ (Artikel 4).

III. Da $a^m a^{m+1}=D+b^m b^m$ und $a^{m+1}$ nicht $<a^m$ ist, ist $a^m a^m$ nicht $>D+b^m b^m$. Da $b^m$ nicht $>\frac12 a^m$ ist, ist $a^m a^m$ nicht $>D+\frac14 a^m a^m$ und $\frac34 a^m a^m$ nicht $>D$; schließlich ist $a^m$ \emph{nicht} $>\surd{\frac43 D}$.

\emph{Beispiel.} Die gegebene Form sei $(304,217,155)$, ihre Determinante $=-31$. Hier ergibt sich die Folge der Formen:
\[
(304,217,155),\ (155,-62,25),\ (25,12,7),\ (7,2,5),\ (5,-2,7).
\]
Die letzte ist die gesuchte. --- Ebenso findet man zur Form $(121,49,20)$ mit Determinante $=-19$ die äquivalenten Formen $(20,-9,5)$, $(5,-1,4)$, $(4,1,5)$; somit ist $(4,1,5)$ die gesuchte Form.

Solche Formen $(A,B,C)$ mit negativer Determinante, bei denen $A$ weder größer als $\surd{\frac43 D},C$ noch kleiner als $2B$ ist, nennen wir \emph{reduzierte Formen}. Zu jeder Form mit negativer Determinante lässt sich also eine eigentlich äquivalente reduzierte Form finden.

\Needspace{8\baselineskip}
\Article{172}
\textsc{Aufgabe.} \emph{Die Bedingungen finden, unter denen zwei nicht identische reduzierte Formen mit derselben Determinante $-D$, $(a,b,c)$, $(a',b',c')$, eigentlich äquivalent sein können.}

\emph{Lösung.} Wir dürfen annehmen, dass $a'$ nicht $>a$ ist und die Form $axx+2bxy+cyy$ in $a'x'x'+2b'x'y'+c'y'y'$ übergeht, und zwar durch die eigentliche Substitution $x=\alpha x'+\beta y'$, $y=\gamma x'+\delta y'$. Dann gelten die Gleichungen
\[
\begin{aligned}
a\alpha\alpha+2b\alpha\gamma+c\gamma\gamma&=a'&\quad\mathinner{.\,.\,.\,.\,.\,.\,.\,.}\quad&[1]\\
a\alpha\beta+b(\alpha\delta+\beta\gamma)+c\gamma\delta&=b'&\quad\mathinner{.\,.\,.\,.\,.\,.\,.}\quad&[2]\\
\alpha\delta-\beta\gamma&=1&\quad\mathinner{.\,.\,.\,.\,.\,.\,.\,.\,.\,.\,.}\quad&[3]
\end{aligned}
\]
Aus [1] folgt $aa'=(a\alpha+b\gamma)^2+D\gamma\gamma$; somit ist $aa'$ positiv. Da $ac=D+bb$, $a'c'=D+b'b'$, sind auch $ac$, $a'c'$ positiv; also haben $a,a',c,c'$ alle dasselbe Vorzeichen. Aber sowohl $a$ als auch $a'$ sind nicht $>\surd{\frac43 D}$, also ist $aa'$ nicht $>\surd{\frac43 D}$. Umso weniger kann daher $D\gamma\gamma\ (=aa'-(a\alpha+b\gamma)^2)$ größer als $\frac43 D$ sein. Somit ist $\gamma$ entweder $=0$ oder $=\pm1$.

I. Ist $\gamma=0$, so folgt aus [3] entweder $\alpha=1.\,\delta=1$ oder $\alpha=-1$, $\delta=-1$.
\SourceLeaf{153}
In beiden Fällen ergibt sich aus [1] $a'=a$ und aus [2] $b'-b=\pm\beta a$. Aber $b$ ist nicht $>\frac12 a$, und $b'$ ist nicht $>\frac12 a'$, also auch nicht $>\frac12 a$. Daher kann die Gleichung $b'-b=\pm\beta a$ nur bestehen, wenn

\emph{entweder} $b=b'$ gilt; daraus würde $c'=\frac{b'b'+D}{a'}=\frac{bb+D}{a}=c$ folgen, und die Formen $(a,b,c)$, $(a',b',c')$ wären entgegen der Voraussetzung identisch;

\emph{oder} $b=-b'=\pm\frac12 a$ gilt. Auch in diesem Fall ist $c'=c$, und die Form $(a',b',c')$ ist $(a,-b,c)$, \emph{das heißt} die zur Form $(a,b,c)$ entgegengesetzte. Zugleich ist klar, dass diese Formen wegen $2b=\pm a$ ambig sind.

II. Ist $\gamma=\pm1$, so ergibt sich aus [1] $a\alpha\alpha+c-a'=\pm2b\alpha$. Aber $c$ ist nicht kleiner als $a$, also auch nicht kleiner als $a'$. Daher ist $a\alpha\alpha+c-a'$ oder $2b\alpha$ sicher nicht kleiner als $a\alpha\alpha$. Da $2b$ nicht größer als $a$ ist, ist $\alpha$ nicht kleiner als $\alpha\alpha$; somit gilt notwendig entweder $\alpha=0$ oder $=\pm1$.

1) Ist $\alpha=0$, so folgt aus [1] $a'=c$. Da $a$ weder größer als $c$ noch kleiner als $a'$ ist, gilt notwendig $a'=a=c$. Weiter ergibt sich aus [3] $\beta\gamma=-1$ und damit aus [2] $b+b'=\pm\delta c=\pm\delta a$. Hieraus folgt wie in (I), dass

\emph{entweder} $b=b'$ gilt; dann wären die Formen $(a,b,c)$, $(a',b',c')$ entgegen der Voraussetzung identisch;

\emph{oder} $b=-b'$ gilt; dann sind die Formen $(a,b,c)$, $(a',b',c')$ entgegengesetzt.

2) Ist $\alpha=\pm1$, so folgt aus [1] $\pm2b=a+c-a'$. Da nun $a$ nicht kleiner ist und auch $c<a'$ nicht gilt, ist $2b$ nicht $<a$ und nicht $<c$. Aber $2b$ ist auch nicht $>a$ und nicht $>c$. Daher gilt notwendig $\pm2b=a=c$ und somit nach der Gleichung $\pm2b=a+c-a'$ auch $=a'$. Aus [2] ergibt sich also
\[
b'=a(\alpha\beta+\gamma\delta)+b(\alpha\delta+\beta\gamma)
\]
oder wegen $\alpha\delta-\beta\gamma=1$
\[
b'-b=a(\alpha\beta+\gamma\delta)+2b\beta\gamma=a(\alpha\beta+\gamma\delta\pm\beta\gamma)
\]
also notwendig, wie zuvor,

\emph{entweder} $b=b'$, sodass die Formen $(a,b,c)$, $(a',b',c')$ entgegen der Voraussetzung identisch wären;

\emph{oder} $b=-b'$, sodass jene Formen entgegengesetzt sind. Zugleich sind die Formen in diesem Fall wegen $a=\pm2b$ ambig.

Aus allem zusammen folgt, dass die Formen $(a,b,c)$, $(a',b',c')$ nur dann eigentlich äquivalent sein können, wenn sie entgegengesetzt und zugleich \emph{entweder} ambig sind \emph{oder} $a=c=a'=c'$ gilt.
\SourceLeaf{154}
Dass die Formen $(a,b,c)$, $(a',b',c')$ in diesen Fällen eigentlich äquivalent sind, hätte man auch leicht im Voraus sehen können. Denn entgegengesetzte Formen müssen uneigentlich und, wenn sie außerdem ambig sind, auch eigentlich äquivalent sein. Ist dagegen $a=c$, so ist die Form $(\frac{D+(a-b)^2}{a},a-b,a)$ der Form $(a,b,c)$ benachbart und damit äquivalent. Wegen $D+bb=ac=aa$ gilt aber $\frac{D+(a-b)^2}{a}=2a-2b$, und die Form $(2a-2b,a-b,a)$ ist ambig. Also ist $(a,b,c)$ auch ihrer entgegengesetzten Form eigentlich äquivalent.

Ebenso leicht lässt sich nun entscheiden, wann zwei reduzierte Formen $(a,b,c)$, $(a',b',c')$, die nicht entgegengesetzt sind, uneigentlich äquivalent sein können. Sie sind nämlich uneigentlich äquivalent, wenn $(a,b,c)$, $(a',-b',c')$, die nicht identisch sind, eigentlich äquivalent sind, und umgekehrt. Damit ist die Bedingung, unter der \emph{jene} uneigentlich äquivalent sind, dass sie identisch und außerdem entweder ambig sind oder $a=c$ gilt. --- Reduzierte Formen, die weder identisch noch entgegengesetzt sind, können weder eigentlich noch uneigentlich äquivalent sein.

\Needspace{8\baselineskip}
\Article{173}
\textsc{Aufgabe.} \emph{Für zwei gegebene Formen $F$ und $F'$ mit derselben negativen Determinante feststellen, ob sie äquivalent sind.}

\emph{Lösung.} Man sucht zwei reduzierte Formen $f$, $f'$, die den Formen $F$, $F'$ jeweils eigentlich äquivalent sind. Sind die Formen $f$, $f'$ eigentlich, uneigentlich oder auf beide Arten äquivalent, so gilt das auch für $F$, $F'$. Sind dagegen $f$, $f'$ auf keine Art äquivalent, so auch $F$, $F'$ nicht.

Nach dem vorigen Artikel gibt es vier mögliche Fälle:

1) Sind $f$, $f'$ weder identisch noch entgegengesetzt, so sind $F$, $F'$ auf keine Art äquivalent.

2) Sind $f$, $f'$ \emph{erstens} entweder identisch oder entgegengesetzt und \emph{zweitens} entweder ambig oder haben gleiche äußere Terme, so sind $F$, $F'$ sowohl eigentlich als auch uneigentlich äquivalent.

3) Sind $f$, $f'$ identisch, aber weder ambig noch mit gleichen äußeren Termen, so sind $F,F'$ nur eigentlich äquivalent.

4) Sind $f$, $f'$ entgegengesetzt, aber weder ambig noch mit gleichen äußeren Termen, so sind $F,F'$ nur uneigentlich äquivalent.

\emph{Beispiel.} Zu den Formen $(41,35,30)$, $(7,18,47)$ mit Determinante $=-5$ findet man die äquivalenten reduzierten Formen $(1,0,5)$, $(2,1,3)$; also sind jene auf keine Art äquivalent. --- Zu den Formen $(23,38,63)$, $(15,20,27)$ ist dagegen dieselbe reduzierte
\SourceLeaf{155}
Form $(2,1,3)$ äquivalent. Da diese zugleich ambig ist, sind die Formen $(23,38,63)$, $(15,20,27)$ sowohl eigentlich als auch uneigentlich äquivalent. --- Zu den Formen $(37,53,78)$, $(53,73,102)$ sind die reduzierten Formen $(9,2,9)$, $(9,-2,9)$ äquivalent. Da diese entgegengesetzt sind und ihre äußeren Terme gleich sind, sind die gegebenen Formen sowohl eigentlich als auch uneigentlich äquivalent.

\Needspace{8\baselineskip}
\Article{174}
Die Anzahl aller reduzierten Formen mit einer gegebenen Determinante $-D$ ist immer endlich und im Verhältnis zur Zahl $D$ recht klein. Die Formen selbst lassen sich auf zwei Arten finden. Wir bezeichnen die reduzierten Formen der Determinante $-D$ allgemein mit $(a,b,c)$; es sind also alle Werte von $a,b,c$ zu bestimmen.

\emph{Erste Methode.} Für $a$ nimmt man alle positiven und negativen Zahlen, die nicht größer als $\surd{\frac43 D}$ sind und für die $-D$ ein quadratischer Rest ist. Für jedes $a$ setzt man $b$ der Reihe nach gleich allen Werten des Ausdrucks $\surd{-D}\ (\mathrm{mod.}\ a)$, die nicht größer als $\frac12 a$ sind, jeweils positiv und negativ genommen. Für $c$ setzt man bei jedem bestimmten Wertepaar von $a,b$ den Wert $=\frac{D+bb}{a}$. Entstehen dabei Formen, in denen $c<a$ gilt, so verwirft man sie; die übrigen sind offenbar reduziert.

\emph{Zweite Methode.} Für $b$ nimmt man alle positiven und negativen Zahlen, die nicht größer als $\frac12\surd{\frac43 D}$ oder $\surd{\frac13 D}$ sind. Für jedes $b$ zerlegt man $bb+D$ auf alle möglichen Arten in zwei Faktoren (auch unter Berücksichtigung verschiedener Vorzeichen), die beide nicht kleiner als $2b$ sind. Den einen Faktor, und zwar bei ungleichen Faktoren den kleineren, setzt man $=a$, den anderen $=c$. Da $a$ nicht $>\surd{\frac43 D}$ ist, sind alle so entstehenden Formen offenbar reduziert. --- Schließlich ist klar, dass jede reduzierte Form mit jeder der beiden Methoden gefunden wird.

\emph{Beispiel.} Sei $D=85$. Die Grenze für die Werte von $a$ ist hier $\surd{\frac{340}{3}}$, die zwischen $10$ und $11$ liegt. Die Zahlen zwischen $1$ und $10$ (einschließlich), für die $-85$ ein Rest ist, sind $1,2,5,10$. Daraus ergeben sich zwölf Formen:
\[
\begin{gathered}
(1,0,85),\ (2,1,43),\ (2,-1,43),\ (5,0,17),\\
(10,5,11),\ (10.\,-5,11);\ (-1,0,-85),\ (-2,1,-43),\\
(-2,-1,-43).\ (-5,0,-17),\ (-10,5,-11,\ (-10,-5,-11).
\end{gathered}
\]
Bei der anderen Methode ergibt sich als Grenze der Werte von $b$ die Zahl $\surd{\frac{85}{3}}$, die zwischen $5$ und $6$ liegt. Für $b=0$ entstehen die Formen
\SourceLeaf{156}
\[
(1,0,85),\ (-1,0,-85),\ (5,0,17),\ (-5,0,-17),
\]
und für $b=\pm1$ diese: $(2,\pm1,43)$, $(-2,\pm1,-43)$.

Für $b=\pm2$ gibt es keine, da sich $89$ nicht in zwei Faktoren zerlegen lässt, die beide nicht $<4$ sind. Dasselbe gilt für $\pm3$, $\pm4$. Schließlich ergeben sich für $b=\pm5$
\[
(10,\pm5,11),\ (-10,\pm5.\,-11).
\]

\Needspace{8\baselineskip}
\Article{175}
Verwirft man aus allen reduzierten Formen einer gegebenen Determinante von je zwei Formen, die zwar nicht identisch, aber eigentlich äquivalent sind, eine der beiden, so haben die verbleibenden Formen folgende bemerkenswerte Eigenschaft: Jede Form mit derselben Determinante ist einer von ihnen eigentlich äquivalent, und zwar nur einer einzigen (sonst wären einige von ihnen eigentlich äquivalent). Also \emph{lassen sich alle Formen mit derselben Determinante in ebenso viele Klassen einteilen, wie Formen übrig bleiben}, indem man die Formen, die derselben reduzierten Form eigentlich äquivalent sind, derselben Klasse zuordnet. So bleiben für $D=85$ die Formen
\[
\begin{gathered}
(1,0,85),\ (2,1,43),\ (5,0,17),\ (10,5,11),\\
(-1,0.\,-84),\ (-2,1,-43).\ (-5,0,-17),\ (-10,5,-11)
\end{gathered}
\]
übrig. Alle Formen mit Determinante $-85$ lassen sich also in acht Klassen einteilen, je nachdem, ob sie der ersten, der zweiten usw. Form eigentlich äquivalent sind. Offenbar sind Formen derselben Klasse eigentlich äquivalent, während Formen aus verschiedenen Klassen nicht eigentlich äquivalent sein können. Die Klassifikation der Formen werden wir weiter unten viel ausführlicher behandeln. Hier fügen wir nur eine Beobachtung hinzu. Wir haben oben schon gezeigt: Ist die Determinante der Form $(a,b,c)$ negativ $=-D$, so haben $a$ und $c$ dasselbe Vorzeichen (denn es gilt $ac=bb+D$, und dieser Wert ist positiv). Ebenso erkennt man leicht, dass bei äquivalenten Formen $(a,b,c)$, $(a',b',c')$ alle $a,c,a',c'$ dasselbe Vorzeichen haben. Geht nämlich die erste durch die Substitution $x=\alpha x'+\beta y'$, $y=\gamma x'+\delta y'$ in die zweite über, so gilt $a\alpha\alpha+2b\alpha\gamma+c\gamma\gamma=a'$ und daher $aa'=(a\alpha+b\beta)^2+D\gamma\gamma$, was sicher nicht negativ ist. Da aber $a$ nicht null sein kann und $a'=0$ ausgeschlossen ist, ist $aa'$ positiv, und damit haben $a,a'$ dasselbe Vorzeichen.

Damit ist klar, dass die Formen mit positiven äußeren Termen von denen mit negativen äußeren Termen vollständig getrennt sind. Es reicht aus, von den reduzierten Formen nur die mit positiven äußeren Termen zu betrachten, denn
\SourceLeaf{157}
die übrigen sind gleich zahlreich und entstehen aus ihnen, indem man den äußeren Termen die entgegengesetzten Vorzeichen gibt. Dasselbe gilt für die Formen, die aus den reduzierten Formen zu verwerfen oder beizubehalten sind.

\Needspace{25\baselineskip}
\Article{176}
Hier steht nun für einige negative Determinanten eine Tabelle von Formen, nach denen sich alle übrigen Formen mit derselben Determinante in Klassen einteilen lassen. Nach der Bemerkung im vorigen Artikel geben wir nur die Hälfte an, nämlich die Formen mit positiven äußeren Termen.
\[
\begin{array}{r|l}
\multicolumn{1}{r}{D}&\\
1&(1,0,1).\\
2&(1,0,2).\\
3&(1,0,3),\ (2,1,2).\\
4&(1,0,4),\ (2,0,2).\\
5&(1,0,5),\ (2,1,3).\\
6&(1,0,6),\ (2,0,3).\\
7&(1,0,7),\ (2,1,4).\\
8&(1,0,8),\ (2,0,4),\ (3,1,3).\\
9&(1,0,9),\ (2,1,5),\ (3.\,0,3).\\
10&(1,0,10),\ (2,0,5).\\
11&(1,0,11),\ (2,1,6),\ (3,1,4),\ (3,-1,4).\\
12&(1,0,12),\ (2,0,6),\ (3,0,4),\ (4,2,4).
\end{array}
\]
Es wäre unnötig, diese Tabelle hier weiter fortzusetzen, da wir weiter unten zeigen werden, wie sie sich viel zweckmäßiger anordnen lässt.

Also ist jede Form mit Determinante $-1$ der Form $xx+yy$ eigentlich äquivalent, wenn ihre äußeren Terme positiv sind, oder der Form $-xx-yy$, wenn sie negativ sind; jede Form mit Determinante $-2$ und positiven äußeren Termen der Form $xx+2yy$ usw.; jede Form mit Determinante $-11$ und positiven äußeren Termen einer der Formen $xx+11yy$, $2xx+2xy+6yy$, $3xx+2xy+4yy$, $3xx-2xy+4yy$ usw.
